\chapter{Troubleshooting}
\label{ch:trouble}

Building a custom Linux system from source involves many interdependent
components. When something goes wrong, the symptoms may appear far from
the root cause. This chapter documents the most common issues encountered
during Veilbox development and their solutions.

\section{Build Failures}

\subsection{Kernel Compilation Errors}

\paragraph{Symptom: make fails with ``implicit function declaration''}

This typically indicates a GCC version mismatch. The kernel requires GCC
14+ for Linux 7.1.0-rc6.

\begin{lstlisting}[language=bash]
gcc --version   # Verify version
sudo dnf update gcc  # Update if needed
make mrproper   # Clean build tree
\end{lstlisting}

\paragraph{Symptom: ``No rule to make target 'debian/certs/signing\_key.pem'''}

The kernel build system expects signing keys for module signature
verification. If module signing is enabled but no keys are configured:

\begin{lstlisting}[language=bash]
# Disable module signing in kernel config
scripts/config --disable MODULE_SIG
make olddefconfig
\end{lstlisting}

\paragraph{Symptom: ``Kernel panic - not syncing: VFS: Unable to mount root fs''}

The initramfs was not properly embedded. Verify:

\begin{lstlisting}[language=bash]
grep CONFIG_INITRAMFS_SOURCE .config
# Should show: CONFIG_INITRAMFS_SOURCE="initramfs.list"
ls -la initramfs.list  # File must exist
# Check kernel binary contains cpio archive
binwalk -e arch/x86/boot/bzImage | grep "cpio"
\end{lstlisting}

\subsection{Initramfs Issues}

\paragraph{Symptom: ``/init: line 1: syntax error: unexpected word''}

The init script has DOS line endings (CRLF). The kernel's cpio extractor
preserves file contents exactly, so if the editor added \textbackslash r
characters, the shell will choke:

\begin{lstlisting}[language=bash]
# Check for CRLF line endings
file rootfs/etc/init.d/rcS
# Fix if needed: dos2unix rootfs/etc/init.d/rcS
\end{lstlisting}

\paragraph{Symptom: ``/bin/sh: not found''}

BusyBox is missing from the initramfs or the symlink is broken:

\begin{lstlisting}[language=bash]
ls -la rootfs/bin/sh          # Should point to busybox
ls -la rootfs/bin/busybox     # Must exist
file rootfs/bin/busybox       # Should be ELF binary
\end{lstlisting}

\section{Boot Failures}

\subsection{GRUB Boot Failures}

\paragraph{Symptom: ``Geometric error'' or triple fault during GRUB stage 1.5}

The core.img blocklist is corrupted. This happens when the blocklist
is manually patched with incorrect format (see Chapter 6 for the correct
12-byte format). Solution: regenerate core.img with grub2-mkimage and
do NOT patch the blocklist.

\begin{lstlisting}[language=bash]
grub2-mkimage -c embedded.cfg \
    -o core.img -O i386-pc \
    ext2 part_msdos biosdisk search
\end{lstlisting}

\paragraph{Symptom: ``No such partition''}

GRUB cannot find the VEILBOX partition. Common causes:

\begin{itemize}[nosep]
  \item The ext4 filesystem label does not match. Verify with debugfs:
        \texttt{debugfs -R "stats" veilbox.raw} | grep "Filesystem volume name"
  \item The part\_msdos module is missing from core.img. Regenerate with
        \texttt{part\_msdos} in the module list.
  \item The disk image is not MBR-partitioned. Verify with
        \texttt{sfdisk -l veilbox.raw}.
\end{itemize}

\paragraph{Symptom: ``You need to load the kernel first''}

GRUB cannot find the kernel file. Check:

\begin{itemize}[nosep]
  \item The file path in grub.cfg matches the actual location on the
        VEILBOX partition.
  - The ext2 module can read the partition. The VEILBOX partition must
        be ext4 with the ext2 driver (backward compatible).
\end{itemize}

\subsection{QEMU Boot Failures}

\paragraph{Symptom: QEMU exits immediately with ``invalid MZ header''}

The kernel lacks CONFIG\_EFI\_STUB. QEMU's \texttt{-kernel} mode requires
a valid MZ/PE header at the start of the kernel binary:

\begin{lstlisting}[language=bash]
grep CONFIG_EFI_STUB .config
# If not set: echo "CONFIG_EFI_STUB=y" >> .config
\end{lstlisting}

\paragraph{Symptom: Black screen in -nographic mode}

The serial console is not configured:

\begin{lstlisting}[language=bash]
grep CONFIG_SERIAL_8250_CONSOLE .config  # Must be y
# Ensure kernel cmdline includes console=ttyS0
grep CONFIG_CMDLINE .config
\end{lstlisting}

\section{Network Issues}

\paragraph{Symptom: No IP address after boot}

Check the DHCP client:
\begin{lstlisting}
ps | grep udhcpc           # Is it running?
ip addr show eth0          # Does the interface exist?
udhcpc -i eth0 -v          # Run in verbose mode
\end{lstlisting}

If udhcpc reports ``ifconfig failed'', the interface name does not match.
The default script uses \$interface which is set by udhcpc. Verify the
interface name matches in the script.

\section{SSH Issues}

\paragraph{Symptom: ``Permission denied (publickey,password)''}

The Dropbear \texttt{-w} flag was set, which in Dropbear v2025.89+
means ``\textbf{Disallow} root logins'' (the opposite of older versions).
Remove the \texttt{-w} flag from the Dropbear command line.

\paragraph{Symptom: ``Connection refused''}

Dropbear is not running or not listening on the expected port:
\begin{lstlisting}
ps | grep dropbear
netstat -tlnp | grep 22
\end{lstlisting}

\paragraph{Symptom: ``Host key verification failed''}

The guest's host keys changed (e.g., after rebuilding). Remove the old
key from the known\_hosts file:
\begin{lstlisting}[language=bash]
ssh-keygen -R '[localhost]:2222'
\end{lstlisting}

\section{Container Runtime Issues}

\paragraph{Symptom: ``containerd: failed to load cgroup''}

The kernel lacks cgroup support:
\begin{lstlisting}[language=bash]
grep CONFIG_CGROUPS .config           # Must be y
grep CONFIG_MEMCG .config             # Must be y
mount | grep cgroup                   # Must show cgroup mounts
\end{lstlisting}

\paragraph{Symptom: ``runc: permission denied''}

Seccomp is enabled but the kernel lacks seccomp support:
\begin{lstlisting}[language=bash]
grep CONFIG_SECCOMP .config           # Must be y
grep SECCOMP /proc/self/status        # Must show seccomp
\end{lstlisting}

\section{Storage Issues}

\paragraph{Symptom: ``No space left on device''}

The 128 MB state partition is full. Clean up:

\begin{lstlisting}
# Check usage
df -h /mnt/state

# Remove unused container images
nerdctl images
nerdctl rmi <image-id>

# Remove stopped containers
nerdctl rm $(nerdctl ps -aq)

# Clean containerd content store
du -sh /mnt/state/containerd/*
rm -rf /mnt/state/containerd/tmp/*
\end{lstlisting}

\paragraph{Symptom: ``Failed to mount /mnt/state''}

The state partition device is not available. The fallback behavior
mounts a tmpfs, so /mnt/state will exist but will be empty. To debug:

\begin{lstlisting}
# Check available block devices
ls -la /dev/vd* /dev/sd*

# Check kernel messages for storage errors
dmesg | grep -i "vdb\|sdb\|ext4"

# Manually mount
mount /dev/vdb /mnt/state 2>&1
\end{lstlisting}

\section{Init Process Issues}

\paragraph{Symptom: ``Kernel panic - not syncing: Attempted to kill init!''}

The init process (PID 1) crashed. This is fatal because the kernel
has no process to manage. Common causes:

\begin{itemize}[nosep]
  \item \textbf{BusyBox missing}: /sbin/init does not exist. Verify
        the BusyBox binary and symlinks in the initramfs.
  \item \textbf{/etc/inittab missing or malformed}: BusyBox init reads
        inittab on startup. If the file is missing or has syntax errors,
        init may not spawn any processes.
  \item \textbf{Wrong init path}: The kernel command line specifies an
        init path (\texttt{init=/sbin/init}) that does not exist on the
        root filesystem.
  \item \textbf{Shared library missing}: If BusyBox is dynamically linked
        (it should be static in Veilbox), missing libraries prevent it
        from executing.
\end{itemize}

\paragraph{Symptom: ``init: can't run /etc/init.d/rcS''}

The rcS script is missing, has incorrect permissions, or has a shebang
pointing to a non-existent interpreter:

\begin{lstlisting}
# Check permissions
ls -la /etc/init.d/rcS
# Should be: -rwxr-xr-x

# Check shebang
head -1 /etc/init.d/rcS
# Should be: #!/bin/sh

# Verify /bin/sh exists
ls -la /bin/sh
\end{lstlisting}

\section{Containerd Troubleshooting}

\paragraph{Symptom: ``containerd: failed to start daemon: OCI runtime create failed''}

The runc binary is missing or incompatible:

\begin{lstlisting}
# Verify runc exists and is executable
ls -la /usr/bin/runc
file /usr/bin/runc

# Check runc version
runc --version

# Verify containerd can find runc
containerd config default | grep runtime
\end{lstlisting}

\paragraph{Symptom: ``Failed to pull image: dial tcp: lookup registry-1.docker.io''}

DNS resolution is not working:

\begin{lstlisting}
# Check DNS configuration
cat /etc/resolv.conf

# Test DNS resolution
nslookup registry-1.docker.io
# or: ping -c 1 registry-1.docker.io

# Manually set DNS
echo "nameserver 8.8.8.8" > /etc/resolv.conf
echo "nameserver 1.1.1.1" >> /etc/resolv.conf
\end{lstlisting}

\section{Serial Console Issues}

\paragraph{Symptom: No output on serial console}

The serial port is not configured in the kernel or QEMU:

\begin{lstlisting}[language=bash]
# Verify kernel config
grep CONFIG_SERIAL_8250_CONSOLE .config
grep CONFIG_SERIAL_8250 .config

# Verify QEMU serial port
# Should have: -nographic or -serial stdio

# Check kernel command line
cat /proc/cmdline
# Should include: console=ttyS0
\end{lstlisting}
